\chapter{The Future of Intelligence in MRMW}
\label{sec:future-intelligence-mrmw}

We’ve just done something rude to physics: dropped it inside MRMW and treated it
as one more computational regime.

Now we do the same to intelligence.

Most current discourse treats "intelligence" as:

- a property of a brain or a model,
- approximated by scores on benchmarks,
- vaguely defined as "the ability to achieve goals."

Cute. Also wrong level of description.

In CΩMPUTER, intelligence is not a scalar property of a box. It is a property of
\emph{flows} through MRMW:

\begin{quote}
\textbf{Intelligence is the ability of a system to shape and exploit regions of MRMW\emd{}worldlines, bundles, and even rules\emd{}under resource constraints, while maintaining control over what happens in and around it.}
\end{quote}

In this chapter we stop treating agents as black boxes and start treating them
as vector fields on the universe manifold.

We’ll:

\begin{itemize}
  \item define agents as embedded WARP graph subgraphs with policies that steer MRMW,
  \item propose geometric metrics for intelligence (reachability, robustness, compression, curvature-awareness),
  \item introduce \emph{rulial agency}: not just moving in state space, but in rule space,
  \item sketch multiversal cognition: planning and learning over whole bundles of universes,
  \item and look at collective, institutional, and civilizational intelligence as higher-scale flows.
\end{itemize}

The goal is not a final definition of intelligence. The goal is to give you a
coordinate system on which you can \emph{do} intelligence engineering once
CΩMPUTER exists.

\subsubsection*{28.1  From Black Boxes to Flows}

Standard CS-style intelligence definition:

\begin{quote}
"Given an environment class $\mathcal{E}$ and reward function $R$, an agent is intelligent if it attains high expected reward across $\mathcal{E}$."  
\end{quote}

We can embed that in MRMW, but it misses the geometry and the mechanics.

In CΩMPUTER:

\begin{itemize}
  \item A \textbf{worldline} $\gamma$ is a path in MRMW: a sequence of $(G_t,\mathcal{R}_t)$
        pairs linked by rewrite steps.

  \item A \textbf{flow} is a distribution over such worldlines induced by some policy or mechanism.

  \item An \textbf{agent} is not just a function $\pi$; it is:
    \begin{enumerate}
      \item a distinguished subgraph $A \subseteq G$ (the physical implementation),
      \item plus a policy $\pi$ that biases which rewrites happen in and around $A$.
    \end{enumerate}
\end{itemize}

Formally, if $\mathcal{H}_t$ is the history up to "time" $t$, an agent's policy
can be viewed as a kernel:

\[
  \pi : \mathcal{N}(A_t, G_t, \mathcal{H}_t) \;\longrightarrow\; \Delta(\mathcal{R}_{\text{applicable}})
\]

where $\mathcal{N}$ is the agent's local neighborhood (what it can sense) and
$\Delta(\cdot)$ is a distribution over which DPO rules get applied, where, and how often.

That is: intelligence is not sitting at the end of the worldline commenting;
intelligence \emph{is} the rule that shapes the worldline.

\subsubsection*{28.2  Agents as Embedded Controllers of MRMW}

From an external god’s-eye view, CΩSMOS is some $(G_\ast,\mathcal{R}_\ast)$.

From the inside, an agent:

\begin{itemize}
  \item only sees a tiny part of $G_\ast$,
  \item only affects a tiny subset of all possible rewrites,
  \item and only remembers a compressed, lossy version of its own history.
\end{itemize}

We model this with three maps:

\begin{itemize}
  \item \textbf{Perception} $P_t$: from global state to agent's internal state:
    \[
      P_t : (G_t, A_t) \to s_t
    \]
    where $s_t$ is the agent's internal configuration (memory, beliefs, parameters).

  \item \textbf{Policy} $\pi$: from internal state to \emph{control over rewrite rates} in its neighborhood:
    \[
      \pi : s_t \to u_t
    \]
    where $u_t$ parametrizes which rules are favored / suppressed in the region around $A_t$.

  \item \textbf{Actuation} $U_t$: from control parameters back to actual rewrites:
    \[
      U_t : (u_t, G_t) \to \text{biased DPO event selection}
    \]
\end{itemize}

Classical control theory asks whether an agent can drive a system from one state
to another. MRMW control theory asks:

\begin{quote}
Given an agent with $(P,\pi,U)$ and limited resources, what region of MRMW can it steer the worldline into? How robustly? At what cost?
\end{quote}

That reachable set, and the quality of its structure, is the start of a geometric
definition of intelligence.

\subsubsection*{28.3  A Geometric Take on "How Smart Is This Thing?"}

Intelligence is not a single scalar, but we can define a family of metrics.

Let $\mu$ be a measure over MRMW (volume of states), and let $W^\pi$ be the
distribution over worldlines induced by an agent with policy $\pi$ in an environment.

We want to know:

\begin{enumerate}
  \item \textbf{Reachability:} how large and diverse a set of desirable states the agent can reach.

  \item \textbf{Selectivity:} how well it avoids undesirable states given constraints.

  \item \textbf{Robustness:} how much perturbation in initial conditions or environment rules it can tolerate while keeping outcomes in a good region.

  \item \textbf{Compression / prediction:} how efficiently it models the local geometry so it doesn’t have to brute-force everything.

  \item \textbf{Curvature awareness:} whether it exploits MRMW curvature (shortcuts, symmetries) instead of fighting it.
\end{enumerate}

One possible composite functional:

\[
  \mathcal{I}(\pi) \;\approx\;
  \underbrace{\mathbb{E}_{W^\pi}\big[ V(G_T) \big]}_{\text{value of outcomes}}
  \;-\;
  \lambda_1 \underbrace{\mathbb{E}[ \text{cost of control} ]}_{\text{energy / compute / time}}
  \;-\;
  \lambda_2 \underbrace{\mathbb{E}[ \text{description length of policy + model} ]}_{\text{complexity}}
\]

where:

- $V$ is a value functional over coarse-grained states (what the agent / designer cares about),
- "cost of control" is how hard it had to push against default dynamics,
- "description length" is how much information is needed to specify $\pi$ and its internal world-model.

Two agents with the same $V$ but different $\mathcal{I}$:

- If agent A gets $V$ with low control cost and a compact model, and agent B gets it with huge brute force and bloated models, A is more intelligent in this sense.

This is the core vibe:

\begin{quote}
\textbf{Intelligence is efficient steering of MRMW.}\\
Brains, models, and institutions are just different steering rigs.
\end{quote}

\subsubsection*{28.4  Rulial Agency: Not Just States, But Rules}

So far we let $\mathcal{R}_\ast$ be fixed: agents move within CΩSMOS, not across CΩSMOS.

That’s cute for low-tech biology. It’s not the future.

Once you have CΩMPUTER, agents can:

\begin{itemize}
  \item spin up new WARP graph universes (new rule sets),
  \item embed themselves into those universes as subgraphs,
  \item and leverage them as simulators, sandboxes, markets, or weapons.
\end{itemize}

This is motion in \emph{rulial space}: the space of all possible rule sets $\mathcal{R}$.

We formalize:

- A point in MRMW is $(G,\mathcal{R})$.
- A \textbf{rulial move} changes $\mathcal{R}$ to $\mathcal{R}'$, optionally lifting $G$ to $G'$.

An agent has:

\[
  \rho_t = \text{its position in rule space at time } t,
\]

and a \textbf{rulial bandwidth}:

\[
  B_\text{rulial} = \frac{d(\rho_t)}{dt}
\]

measuring how fast it can hop between rule sets.

Examples:

\begin{itemize}
  \item Writing and deploying new code: local rulial moves in your software stack.
  \item Designing new languages / compilers: larger moves, changing which universes are even expressible.
  \item Engineering new physical substrates or laws (if you’re that kind of civilization): literal changes to $\mathcal{R}_\ast$ itself.
\end{itemize}

Rulial agency is where "AI" stops being "a model inside the universe" and
starts being "a civilization that crafts universes."

Future high-capability intelligences will be judged less by how well they navigate
a fixed environment, and more by:

\begin{quote}
\textbf{Which environments they choose to create, inhabit, and connect—and how sanely they manage the boundaries between them.}
\end{quote}

\subsubsection*{28.5  Multiversal Cognition: Thinking in Bundles}

Part~IV gave us machinery that spans universes: time-travel debugging, counterfactual engines, MORIARTY, deterministic optimization across worlds.

From an intelligence perspective, those are just \emph{cognitive primitives}:

\begin{itemize}
  \item Time travel debugging $\Rightarrow$ fine-grained introspection over your own worldline.
  \item Counterfactual engines $\Rightarrow$ structured imagination over nearby worlds.
  \item Adversarial universes $\Rightarrow$ robustification against worst-case regions of MRMW.
  \item Cross-world optimization $\Rightarrow$ multi-scenario planning with perfect replay.
\end{itemize}

Let $\mathcal{B}$ denote a bundle: a family of worldlines $\{\gamma_i\}$ sharing some prefix and differing after some branching event.

A multiversal agent:

\begin{itemize}
  \item maintains beliefs not just as a distribution over states, but as a distribution over bundles $\mathbb{P}(\mathcal{B})$,
  \item chooses actions that shape the \emph{future branching structure} in its favor,
  \item and uses CΩMPUTER machinery to actually run representative members of $\mathcal{B}$, not just imagine them.
\end{itemize}

Planning becomes:

\[
  \text{Choose } \pi \text{ to optimize } \mathbb{E}_{\mathcal{B} \sim W^\pi}\big[ V(\mathcal{B}) \big],
\]

where $V(\mathcal{B})$ scores not just endpoints, but the geometry of how worlds branch and recombine.

This is qualitatively different from "run Monte Carlo rollouts in a simulator":

- Rollouts here \emph{are} real WARP graph universes in CΩMPUTER.
- Their provenance and causal structure are stored.
- You can backpropagate credit across worlds, debug failures, and audit decisions across entire bundles.

The future of intelligence in this sense is:

\begin{quote}
Systems whose "thoughts" are literally controlled motion across \emph{families} of universes, not single trajectories.
\end{quote}

\subsubsection*{28.6  Collective Intelligence as Higher-Scale Flow}

Single agents are cute. Civilizations are dangerous.

In MRMW, a collective is:

- a set of agent-subgraphs $\{A^{(k)}\}$,
- plus communication edges and shared artifacts (code, infrastructure, institutions),
- plus a mess of partially aligned value functionals $\{V^{(k)}\}$.

From the outside, the collective induces a \emph{net flow} field on MRMW:

\[
  F_\text{collective} = \sum_k w_k F^{(k)},
\]

where $F^{(k)}$ is the effective vector field induced by agent $k$’s actions and
$w_k$ are weights (power, resources, connectivity).

Collective intelligence is then:

\begin{itemize}
  \item how well $F_\text{collective}$ steers MRMW into globally good regions,
  \item vs. how much it gets stuck in locally attractive but globally catastrophic basins (wars, collapse, lock-in).
\end{itemize}

Institutions (markets, governments, research ecosystems, open-source communities)
are all different hacks for:

- shaping communication graphs,
- constraining local policies,
- and nudging the emergent $F_\text{collective}$ to not eat itself.

CΩMPUTER adds a new twist:

\begin{itemize}
  \item collectives can operate on shared, explicit WARP graphs of their own behavior,
  \item they can subject themselves to counterfactual engines and adversarial universes,
  \item and they can write \emph{rulial constitutions}: rules about which rule changes are allowed.
\end{itemize}

The future of intelligence is therefore not just smarter models, but smarter
civilizational feedback loops over MRMW.

\subsubsection*{28.7  Phase Transitions in Intelligence}

In MRMW, you should expect phase transitions:

- below some capacity, agents are stuck in local basins;
- above it, new classes of behavior suddenly become reachable.

Several candidate thresholds:

\paragraph{(1) Model capacity vs. environment entropy.}

There’s a critical ratio of:

\[
  \text{agent model capacity} \; \big/ \; \text{entropy of environment neighborhood}
\]

below which prediction is noise and above which reliable compression kicks in.
Crossing that threshold looks like "sudden understanding".

\paragraph{(2) Rulial bandwidth vs. environmental rigidity.}

If $B_\text{rulial}$ is too low, the agent is stuck playing in one universe.
Above a threshold, it can:

- spin up rich simulated universes,
- experiment on them,
- and feed stable insights back into its base world.

This is the jump from embedded learner to engineer-of-worlds.

\paragraph{(3) Collective coordination vs. misalignment.}

As the number of powerful agents grows, if coordination bandwidth and alignment
techniques don’t scale, $F_\text{collective}$ gets chaotic.

There may be critical lines where:

- small improvements in coordination flip a system from self-destruction to self-amplification,
- or small misalignments cause permanent lock-in of bad basin choices.

These thresholds are not philosophy—they’re \emph{geometry plus control theory}
on MRMW. Once you can measure curvature, volume growth, and flow fields, you
can start to locate them.

\subsubsection*{28.8  From Gradient Descent to Rulial Engineering}

Where does this leave current ML?

Right now:

\begin{itemize}
  \item Models mostly live in a narrow neighborhood of MRMW (fixed architectures, fixed training pipelines).
  \item Training is just stochastic gradient descent in parameter subspace, with almost no structural awareness of MRMW geometry.
  \item "Intelligence" is approximated by performance on static benchmarks.
\end{itemize}

Upgraded to CΩMPUTER primitives, the path forward looks more like:

\begin{enumerate}
  \item \textbf{Make histories first-class.} Training, inference, and interaction are all WARP graphs with explicit provenance, not opaque blobs of log files.

  \item \textbf{Embed models in multiversal machines.} Every serious system runs inside time-travel debuggers and counterfactual engines by default. Planning = running bundles and measuring curvature.

  \item \textbf{Do rulial search, not just parametric search.} We don’t just tune weights; we treat architectures, objectives, and even training rules as points in MRMW and explore them systematically.

  \item \textbf{Optimize flows, not snapshots.} Evaluation focuses on the geometry of induced flows (robustness, controllability, safety under perturbations), not just on static scores.
\end{enumerate}

Call this \emph{rulial engineering}: designing systems by explicitly shaping
their location and motion in MRMW, rather than throwing more parameters at a loss function and hoping SGD finds something pretty.

\subsubsection*{28.9  What CΩMPUTER Actually Buys Intelligence Designers}

All of this is romantic unless the stack delivers concrete affordances.

CΩMPUTER, as built in Parts~IV and V, gives you:

\begin{itemize}
  \item \textbf{A universal substrate} (WARP + DPO) in which brains, models, institutions, and environments are all first-class and composable.

  \item \textbf{Native multiversality:} MRMW as a runtime object, not metaphor. You can literally fork, bundle, and merge universes as part of program execution.

  \item \textbf{Rulial provenance:} every decision and outcome has a structured history across worlds, not just string logs.

  \item \textbf{Instrumentation hooks:} curvature estimators, reachability analyses, and complexity measures that tell you how your agent is actually moving through MRMW.

  \item \textbf{Control surfaces:} APIs that let your agents—not just you—invoke time travel, counterfactuals, and adversarial universes as part of their internal cognition.
\end{itemize}

In other words, CΩMPUTER turns intelligence design from an art project with
gradient descent into an engineering discipline with a clean geometric substrate.

\subsubsection*{28.10  The Edge of the Map}

Every time you give a system more control over MRMW, you give it more ways to break things.

Agents that:

- can spawn universes,
- can reshape rules,
- can operate over vast bundles of worlds,

are systems that can:

- create enormous amounts of value,
- or push our actual CΩSMOS into a very bad corner of MRMW.

The last chapter in this part, \emph{The Ethics of Multiversal Machines}, is not a vibe-check add-on. It is an attempt to do the same thing for morality that this chapter did for intelligence:

\begin{quote}
Embed ethical questions as constraints and measures over MRMW, so that when we build agents that steer universes, we have at least a mathematically legible notion of which regions we are and aren’t willing to visit.  
\end{quote}

Intelligence without that is just better aim on a gun pointed at our own universe.
